\chapter{Theoretical Background}
\label{ch:theory}

This chapter develops, from first principles, every theoretical tool shmap relies on. No prior
familiarity with sketch-based mapping is assumed.

\section{$k$-mers}

A \textbf{$k$-mer} is a contiguous substring of length $k$ of a DNA sequence. A sequence $S$ of
length $n$ over $\Sigma=\{A,C,G,T\}$ has exactly $n-k+1$ $k$-mers (for $n \ge k$), one starting at
each position $i \in [0, n-k]$: $S[i \mathinner{.\,.} i{+}k)$.

$k$-mers are the atomic unit of comparison in almost all modern sequence analysis: instead of
comparing sequences character-by-character (which is what alignment does, and is expensive),
tools compare the \emph{sets} of $k$-mers two sequences share. Two sequences that share many
$k$-mers are likely to be similar/overlapping/from the same genomic locus; two that share almost
none are likely unrelated. This is the basis of alignment-free sequence comparison.

Two design parameters govern how well $k$-mers work as a similarity proxy:
\begin{itemize}
  \item If $k$ is too small, $k$-mers recur very frequently by chance alone (there are only
    $4^k$ distinct $k$-mers; for $k=8$ that is $65{,}536$, tiny compared to a genome of $10^9$
    bases), so a shared $k$-mer carries little evidence of true homology.
  \item If $k$ is too large, a single sequencing error (common in long reads, particularly Oxford
    Nanopore data, at rates of $1$--$15\%$ depending on chemistry) destroys every $k$-mer that
    overlaps it (up to $k$ of them), so real matches are missed.
\end{itemize}
shmap's default is $k=15$ (flag \code{-k}), with typical experiments in the codebase's own
evaluation harness sweeping $k \in \{17,19,\dots,33\}$.

\subsection{Canonical $k$-mers and strand}
\label{sec:canonical-kmers}

DNA is double-stranded, and a sequencing read may originate from either strand, with no way to
know which \emph{a priori}. A $k$-mer $x$ occurring on the forward strand of the reference is
identical to the \textbf{reverse complement} of $x$ occurring on the reverse strand at the
mirrored position. To make $k$-mer matching strand-agnostic, essentially every tool in this space
computes, for each $k$-mer window, both:
\begin{itemize}[nosep]
  \item the forward-strand hash $h_{fw}$, and
  \item the reverse-complement hash $h_{rc}$ (i.e., the hash of the reverse complement of the
    same window),
\end{itemize}
and uses a canonical combination of the two (commonly $\min(h_{fw}, h_{rc})$, or, as shmap does,
$h_{fw} \oplus h_{rc}$ --- see \S\ref{sec:fracminhash-hash}) as the $k$-mer's identity, plus a
single bit recording which strand ``won'' so that downstream code can recover orientation.

\section{Hashing $k$-mers, and rolling hashes}
\label{sec:rolling-hash}

Because a sequence of length $n$ has $O(n)$ overlapping $k$-mers, naively hashing each one from
scratch costs $O(nk)$. A \textbf{rolling hash} instead updates the previous window's hash in
$O(1)$ to obtain the next window's hash, giving $O(n)$ total. The general recurrence for a
polynomial or XOR-based rolling hash removes the outgoing base's contribution and adds the
incoming base's contribution:
\[
  h_{i+1} = f(h_i, \, \text{out\_base} = S[i], \, \text{in\_base} = S[i+k])
\]
shmap uses an XOR/rotate rolling hash (not a polynomial one); the exact recurrence, including the
reverse-complement side, is derived in \S\ref{sec:fracminhash-hash}.

\section{Sketching: reducing $k$-mer sets to a manageable, comparable size}

Even with an efficient rolling hash, a set of $O(n)$ $k$-mers per sequence is too large to
compare pairwise across, e.g., every read against a whole genome. \textbf{Sketching} techniques
subsample the $k$-mer set down to a small representative subset, in a way that lets similarity
between two sequences be \emph{estimated} from the subsets alone, with a distortion (approximation
error) that shrinks as the subset grows.

\subsection{Minimizer schemes}
\label{sec:minimizers}

The most widely used sketching scheme (used by minimap2, most short-read aligners' seed indices,
Winnowmap, etc.) is the \textbf{minimizer}: slide a window of $w$ consecutive $k$-mers along the
sequence, and keep only the $k$-mer with the smallest hash value in each window (ties broken by
position or a secondary key). This guarantees:
\begin{itemize}[nosep]
  \item \textbf{Bounded density:} at most $2/(w+1)$ of all $k$-mers are ever selected (in
    expectation, for a uniform hash), giving a sketch size proportional to $n/w$.
  \item \textbf{Window guarantee:} any window of $w$ consecutive $k$-mers has \emph{at least one}
    selected minimizer, so no region of the sequence is left completely unsampled --- important
    for finding short matches / handling indels robustly.
\end{itemize}
The window guarantee is minimizers' main strength; their main weakness is that the selected set is
sensitive to single-base edits near a window boundary (a phenomenon studied extensively in the
minimizer/syncmer literature), and that sketch size is governed by $w$, not by a similarity-space
quantity directly.

\subsection{MinHash}
\label{sec:minhash}

An older, different sketching idea (Broder, 1997) underlies Mash and sourmash. Given a hash
function $h$ assumed to distribute $k$-mers uniformly over $[0, H_{\max}]$, the \textbf{MinHash
sketch} of a $k$-mer set $A$ of size $n$ keeps only the $n$ smallest-hash elements for a
fixed target sketch size (\emph{bottom-$n$ sketch}), or equivalently the single minimum hash
value repeated over many independent hash functions (\emph{classic MinHash}). The key theoretical
result is:
\[
  \Pr[\min h(A) = \min h(B)] = \frac{|A \cap B|}{|A \cup B|} = J(A,B)
\]
i.e., the probability that two sets' minimum hashes coincide equals their \textbf{Jaccard
similarity}. Averaging this indicator over $m$ independent hash functions (or, equivalently,
comparing the $m$ smallest of a single hash function's outputs) gives an unbiased estimator of
$J(A,B)$ with standard error $O(1/\sqrt{m})$.

MinHash's practical drawback for genomic sketching: a fixed-size bottom-$n$ sketch of a short
sequence and of a long sequence look the same size, so a short sequence's sketch is a much denser
(and hence unrepresentative, if $n$ is not scaled per-sequence) sample of its $k$-mers. This
motivates \textbf{FracMinHash}.

\subsection{FracMinHash (``Mash Screen'' / scaled MinHash)}
\label{sec:fracminhash-theory}

\textbf{FracMinHash} (used by Mash Screen, sourmash's ``scaled'' sketches, and shmap) fixes not a
sketch \emph{size} but a hash-space \textbf{fraction} $r \in (0,1]$ (shmap's flag \code{-r},
default $0.05$): a $k$-mer with hash $h$ is kept in the sketch iff
\[
  h \le r \cdot H_{\max}
\]
where $H_{\max}$ is the maximum possible hash value (for a 64-bit hash, $H_{\max}=2^{64}-1$).
Because $h$ is (assumed) uniformly distributed, this keeps, in expectation, a fraction $r$ of all
$k$-mers, so \textbf{sketch size scales linearly with sequence length}: $|\text{sketch}(S)|
\approx r\cdot(|S|-k+1)$, for any $|S|$. This is the key advantage over both fixed-size MinHash
and (to a lesser extent) minimizers: one parameter $r$ controls sampling density uniformly across
sequences of any length, and, unlike bottom-$n$ MinHash, two FracMinHash sketches computed with
the \emph{same} threshold are directly comparable set-theoretically (their union/intersection
correspond exactly to the union/intersection of the full $k$-mer sets restricted to the
sub-threshold region) --- no renormalization needed.

\paragraph{Jaccard and containment estimation.} For two sketches $\hat A = A \cap [0, rH_{\max}]$
and $\hat B = B \cap [0, rH_{\max}]$ built with the \emph{same} threshold,
\[
  \widehat{J(A,B)} = \frac{|\hat A \cap \hat B|}{|\hat A \cup \hat B|}
  \qquad\text{estimates}\qquad
  J(A,B)=\frac{|A\cap B|}{|A \cup B|},
\]
and, when $|A| \ll |B|$ (e.g.\ $A$ a read, $B$ the whole genome) the more relevant
\textbf{containment index}
\[
  C(A, B) = \frac{|A \cap B|}{|A|}
\]
(``what fraction of the read's $k$-mers are found somewhere in the reference'') is estimated
directly (without needing to also estimate $|A \cup B|$) as
\[
  \widehat{C(A,B)} = \frac{|\hat A \cap \hat B|}{|\hat A|}.
\]
shmap's two headline metrics, \code{Containment} and \code{Jaccard} (CLI flag \code{-m}), are
exactly these two estimators, computed not over the whole reference but over a small candidate
\emph{bucket} region (Chapter~\ref{ch:scoring}); shmap additionally offers two alternative,
bucket-native metrics (\code{bucket\_SH}, \code{bucket\_LCS}) discussed in
\S\ref{sec:metric-variants}.

\subsection{FracMinHash vs.\ minimizers, concretely}
\label{sec:fracminhash-vs-minimizer}

\begin{center}
\begin{tabular}{@{}p{4cm}p{5.4cm}p{5.4cm}@{}}
\toprule
 & \textbf{Minimizer ($w$-window)} & \textbf{FracMinHash (ratio $r$)} \\
\midrule
Sketch size & $\approx 2n/(w{+}1)$ & $\approx rn$ \\
Coverage guarantee & every $w$-window has $\ge 1$ selected $k$-mer & none: a long low-hash-density
  stretch can (rarely) be unsampled \\
Comparable across lengths & only if $w$ fixed and sequences aligned to the same coordinate frame
  & yes, directly, by construction \\
Used by & minimap2, Winnowmap, most short-read aligners & Mash Screen, sourmash, \textbf{shmap} \\
\bottomrule
\end{tabular}
\end{center}

shmap explicitly favors the containment/Jaccard-estimation cleanliness of FracMinHash over
minimizers' coverage guarantee; its seed-heuristic pruning (\S\ref{sec:seed-heuristic-theory}) is
in fact a direct consequence of working in FracMinHash's clean estimation framework.

\section{Seed-and-extend / seed-chain-extend mapping}
\label{sec:seed-extend}

Given sketches of $P$ (the read) and $T$ (the reference), the generic mapping pipeline is:
\begin{enumerate}
  \item \textbf{Seed:} find every position where a $k$-mer of $P$'s sketch also occurs in $T$'s
    sketch (a \emph{hit}); a hit is a pair (position in $P$, position in $T$).
  \item \textbf{Chain / cluster:} group hits that are roughly co-linear (i.e., consistent with one
    contiguous alignment: as position-in-$P$ increases, position-in-$T$ increases by a
    comparable amount, same strand) into candidate regions.
  \item \textbf{Score / extend:} for each candidate region, compute a similarity score (or, in
    tools that go further, a base-level alignment) and report the best (and, for mapping quality,
    second-best) candidate.
\end{enumerate}
Tools differ enormously in step 2. minimap2 runs an $O(N \log N)$ or $O(N^2)$ (bounded) dynamic
program over all $N$ hits to find the highest-scoring chain of co-linear anchors (\emph{co-linear
chaining}). shmap instead partitions the reference into fixed-geometry, overlapping
\textbf{buckets} (a form of coarse, position-only clustering, described in
Chapter~\ref{ch:bucketing}) and scores whole buckets rather than chaining individual hits --- far
cheaper, at the cost of being less precise about the exact alignment boundary (which shmap does
not need, since it is a mapper, not an aligner).

\section{Co-linear chaining via Longest Increasing / Common Subsequence}
\label{sec:lis-theory}

One of shmap's four scoring metrics, \code{bucket\_LCS} (\S\ref{sec:metric-variants}), scores a
candidate bucket by the length of the \textbf{longest common subsequence} (LCS) between the order
hits occur in $P$ and the order they occur in $T$ within the bucket. Because both sequences here
are permutations of (a subset of) the same multiset of positions-in-$P$, ``LCS'' reduces to
\textbf{longest increasing subsequence} (LIS) of the sequence of $P$-positions visited in
increasing $T$-order: a long increasing run means many hits are mutually co-linear (consistent
with one real alignment on one strand); a long \emph{decreasing} run means many hits are
co-linear on the \emph{other} strand. shmap computes both (forward and reverse) and takes the max
(\S\ref{sec:lcs-implementation}).

\paragraph{Patience-sorting LIS algorithm.} The classic $O(n\log n)$ algorithm for LIS length
(not full reconstruction) processes the sequence left to right, maintaining an array
$\mathit{tails}$ where $\mathit{tails}[i]$ is the smallest possible tail value of any increasing
subsequence of length $i{+}1$ seen so far. For each new value $x$: binary-search $\mathit{tails}$
for the first entry $\ge x$; if none exists, append $x$ (extends the longest run found so far by
one); otherwise overwrite that entry with $x$ (keeps the same achievable lengths but with a
smaller, more extensible tail). The final array length is the LIS length. This is exactly the
algorithm shmap's \code{lcs\_get\_lis} implements (\S\ref{sec:lcs-implementation}).

\section{The seed heuristic: a provable upper bound for cheap pruning}
\label{sec:seed-heuristic-theory}

This is the theoretical core that gives Map-SHmap its name and its main performance advantage
over exhaustively scoring every candidate region.

Consider a read sketch of size $m = |\hat P|$ (its FracMinHash sketch cardinality). Seeds are
consumed \emph{rarest-first} (Chapter~\ref{ch:seeding}) up to some budget $S \le m$ seeds. For any
candidate bucket, define:
\begin{itemize}[nosep]
  \item $\mathit{seeds}$ = number of seed \emph{occurrences} (counting multiplicity in $P$) that
    have been consumed and attributed to this bucket so far (an occurrence is attributed to a
    bucket the moment its position is examined --- whether or not the reference actually has a
    hit there);
  \item $\mathit{matches}$ = number of those occurrences that \emph{did} find a hit landing inside
    this bucket.
\end{itemize}
Necessarily $\mathit{matches} \le \mathit{seeds} \le m$ (a seed occurrence can match at most once
per bucket, and only $\mathit{seeds}\le m$ occurrences from the read have been consumed at all).
Define the \textbf{seed-heuristic score}
\[
  h_{\mathrm{seed}}(m, \mathit{seeds}, \mathit{matches}) \;=\; 1 - \frac{\mathit{seeds} -
  \mathit{matches}}{m}.
\]

\begin{quote}
\textbf{Claim.} $h_{\mathrm{seed}}$ is a non-increasing function of $\mathit{seeds}$ for fixed
$\mathit{matches}$, and an upper bound on the true containment score
$C = \mathit{matches}_{\mathrm{final}} / m$ that this bucket could achieve \emph{even if every
remaining, not-yet-consumed seed occurrence in $P$ also matched inside this bucket}.
\end{quote}

\emph{Why.} At most $m$ total seed occurrences exist; $\mathit{seeds}-\mathit{matches}$ of the
ones examined so far are \emph{known misses} for this bucket (an occurrence was consumed and did
\emph{not} land in this bucket). Even in the best possible case --- every one of the
$m-\mathit{seeds}$ not-yet-examined occurrences turns out to match this bucket too --- the final
match count can be at most $m - (\mathit{seeds}-\mathit{matches})$ (total occurrences minus known
misses), giving final containment score at most
$\bigl(m-(\mathit{seeds}-\mathit{matches})\bigr)/m = h_{\mathrm{seed}}$. Since known misses only
ever accumulate (both $\mathit{seeds}$ and $\mathit{seeds}-\mathit{matches}$ are non-decreasing as
more occurrences are examined), $h_{\mathrm{seed}}$ can only decrease or stay the same as more
seeds are consumed for a fixed bucket --- it is a monotonically tightening upper bound.

\paragraph{Pruning rule.} If, after extending a bucket with more seed occurrences,
$h_{\mathrm{seed}} < \theta$ for the current acceptance threshold $\theta$ (or the current
best-known score, whichever is the active bound), the bucket is \emph{guaranteed} to be unable to
reach $\theta$ no matter how the remaining seeds resolve, and can be discarded \emph{without
computing its actual score}. This is what makes shmap fast: the overwhelming majority of
candidate buckets a repetitive genome generates are eliminated by this $O(1)$-per-check bound
long before the $O(\text{bucket size})$ refinement sweep (Chapter~\ref{ch:scoring}) would ever
run on them. Chapter~\ref{ch:pruning} gives the exact incremental algorithm
(\code{seed\_heuristic\_pass}) that extends a bucket with just enough additional seeds to either
prove it's prunable or exhaust all seeds, at the smallest possible cost.

\section{Mapping quality (MAPQ)}
\label{sec:mapq-theory}

A mapper is rarely certain a single location is correct, especially in repetitive genomic regions
where a read may score similarly against several loci. \textbf{MAPQ} (following the SAM/PAF
convention, an integer typically in $[0,60]$, though PAF's own histrionic convention allows $255$
for ``unavailable'') encodes, on a Phred-like scale, the mapper's confidence that the reported
best location is correct. The general recipe used by essentially all mappers (originating with
MAQ/BWA, and reused with variations by minimap2 and shmap) is:
\begin{enumerate}[nosep]
  \item Find the best-scoring candidate location, score $f_1$.
  \item Find the best-scoring candidate location that is not simply a restatement of the same
    location (typically: does not overlap the first by more than some fraction), score $f_2\le
    f_1$.
  \item If $f_1$ and $f_2$ are close, the read could plausibly have come from either place ---
    report low confidence. If $f_2$ is much worse than $f_1$ (or no second candidate exists at
    all), report high confidence.
\end{enumerate}
shmap's concrete formula, and the specific (somewhat unusual, all-or-nothing rather than smoothly
graduated) way it implements step 3, is given in full in Chapter~\ref{ch:mapq}, including a
discussion of two alternative formulas the original source retains, commented out, as evidence of
in-progress tuning.

\section{Summary of the theoretical pipeline}

Putting the pieces together, shmap's high-level strategy is: \emph{FracMinHash sketch both sides
for length-proportional, directly-comparable $k$-mer subsets; index the reference sketch once;
for each read, seed rarest-$k$-mer-first against the index; accumulate hits into cheap,
overlapping positional buckets; use the seed-heuristic upper bound to discard almost all buckets
for a few arithmetic operations each; refine survivors with an exact sliding-window Jaccard/
containment sweep (or an LCS-based chain score); and report the best two candidates' scores as a
MAPQ estimate.} Chapter~\ref{ch:architecture} now turns this into the concrete system
architecture, and the chapters that follow specify every step precisely enough to implement.
